% =====================================================================
% 第十讲 (Exposé X) 基本群的特殊化理论
% (Théorie de la spécialisation du groupe fondamental)
% A. Grothendieck
% 原书页码: pp. 201-218
% =====================================================================
\chapter{基本群的特殊化理论}

在本讲中，我们限于研究紧合态射中几何纤维的基本群，i.e. 变动的紧合代数概形的基本群. 在以后的讲中，我们将把所用的技巧推广到"在无穷远处"\emph{温和分歧}的平展覆盖. 这将例如给出"三点问题"在阶与特征互素的 Galois 覆盖情形的解 (i.e. 确定 $\mathbf{P}^1_k$ 上至多在给定三点处分歧且在这些点温和分歧的 Galois 覆盖)，以及它的显然变体.

\section{紧合可分态射的同伦正合序列}

\begin{definition}\label{def:X.1.1}
域 $k$ 上的预概形 $X$ 称为可分\index[term]{kefentaishe@可分态射}，或在 $k$ 上可分，若对 $k$ 的任何扩张 $K$，$X \otimes_k K$ 既约. 若 $f \colon X \to Y$ 是预概形态射，称 $f$ 可分，或 $X$ 在 $Y$ 上可分，若 $X$ 在 $Y$ 上平坦，且对每个 $y \in Y$，纤维 $X \otimes_Y k(y)$ 在 $k(y)$ 上可分.
\end{definition}

若 $X$ 是域 $k$ 上的预概形，说它可分也意味着它既约，且诸 $k(x)$ (对 $X$ 的不可约分量的生成点 $x$) 是 $k$ 的可分扩张. 若 $k$ 完备，则说 $X$ 既约与之等价. 注意若 $X$ 在 $Y$ 上可分，则对任何基变换 $Y' \to Y$，$X' = X \times_Y Y'$ 在 $Y'$ 上可分. 在适当的有限性假设下还可证明两个\emph{可分态射}的复合可分. 我们只需下列形式：若 $X$ 在 $Y$ 上可分，且 $X'$ 在 $X$ 上平展，则 $X'$ 在 $Y$ 上可分. 这确实是定义和 I 9.2 的直接推论. 另一方面，"可分态射"假设将借助下列命题起作用：

\begin{proposition}\label{prop:X.1.2}
设 $f \colon X \to Y$ 是紧合可分态射，$Y$ 局部诺特，考虑它的 Stein 分解 $X \xrightarrow{f'} Y' \xrightarrow{f''} Y$ (其中 $f'_*(\mathcal{O}_X) = \mathcal{O}_{Y'}$，$Y'$ 在 $Y$ 上有限且同构于代数 $f_*(\mathcal{O}_X)$ 的谱). 则 $Y'$ 是 $Y$ 的平展覆盖.
\end{proposition}

这个命题将出现在 (EGA III 7)\footnote{参见 EGA III 7.8.10 (i).} 中. 指出证明的原理. 容易归结到 $Y$ 是完备局部环 $A$ 的谱的情形，再对后者作适当的有限平坦扩张 (对应适当的剩余扩张)，可以假设闭点 $y$ 的纤维的连通分量几何连通，这也意味着 $H^0(X_y, \mathcal{O}_{X_y})$ 分解成与 $k = k(y)$ 相同的体的乘积. 此时假设 $X$ 连通 (这是允许的)，将有 $H^0(X_y, \mathcal{O}_{X_y}) = k$，故同态 $A \to H^0(X_y, \mathcal{O}_{X_y})$ 满射. 由此用一般命题 (Künneth 型) 得出 $f_*(\mathcal{O}_X)$ 由 $A$ 上自由模 $B$ 定义，且 $B/\mathfrak{m}B \to H^0(X_y, \mathcal{O}_{X_y}) = k$ 双射. 故此时 $B$ 是 $A$ 上的平展代数，这完成证明.

\begin{theorem}\label{thm:X.1.3}
设 $f \colon X \to Y$ 是紧合可分态射，$Y$ 局部诺特且连通，假设 $f_*(\mathcal{O}_X) = \mathcal{O}_Y$ (这蕴含 $X$ 在 $Y$ 上的诸纤维几何连通，由 \ref{prop:X.1.2} 反之亦然). 设 $y$ 是 $Y$ 的点，$\overline{k(y)}$ 是 $k(y)$ 的代数闭包，$X_{\overline{y}} = X_y \otimes_{k(y)} \overline{k(y)}$. 最后设 $X'$ 是 $X$ 的连通平展覆盖，$\overline{X'}_y = X'_y \otimes_{k(y)} \overline{k(y)}$. 为使存在 $Y$ 的平展覆盖 $Y'$ 和 $X$-同构 $X' \xrightarrow{\sim} X \times_Y Y'$，必要且充分条件是 $\overline{X'}_y$ 在 $X_{\overline{y}}$ 上有截面.
\end{theorem}

令 $Y' = \mathrm{Spec}(h_*(\mathcal{O}_{X'}))$ (其中 $h \colon X' \to Y$ 是复合态射 $X' \to X \to Y$)，只需证明典范 $Y$-态射
\[
X' \longrightarrow X \times_Y Y'
\]
是同构，且 $Y'$ 在 $Y$ 上平展. 由 \ref{prop:X.1.2} 已知 $Y'$ 在 $Y$ 上平展，故 $X \times_Y Y'$ 在 $X$ 上平展，故态射 $X' \to X \times_Y Y'$ 也平展 (\cref{cor:I.4.8}). 而且，$Y'$ 作为连通集 $X'$ 的像连通，故由于 $X$ 在 $Y$ 上有连通纤维 (IX 3.4 和 V 6.9 (iii))，$X \times_Y Y'$ 连通. 故为证 $X' \to X \times_Y Y'$ 是同构，只需看它在 $X \times_Y Y'$ 的某点的投影次数等于 1. 这容易由 $\overline{X'}_y$ 在 $X_{\overline{y}}$ 上有截面的假设得出，或用 IX 6.6，或更简单地注意：只需在基变换 $\mathrm{Spec}(\overline{k(y)}) \to Y$ 后证明 $X \times_Y Y'$ 中存在这样的点，此时显然. 这完成 \ref{thm:X.1.3} 的证明.

考虑到 IX 3.4 和字典 V 6.9 与 V 6.11，可以把 \ref{thm:X.1.3} 写成下列等价形式

\begin{corollary}\label{cor:X.1.4}
用前面关于 $f \colon X \to Y$ 和 $X_{\overline{y}}$ 的记号，设 $\overline{a}$ 是 $X_{\overline{y}}$ 的几何点，$a$ 是它在 $X$ 中的像，$b$ 是它在 $Y$ 中的像. 则下列群同态序列正合：
\[
\pi_1(X_{\overline{y}}, \overline{a}) \longrightarrow \pi_1(X, a) \longrightarrow \pi_1(Y, b) \longrightarrow e.
\]
\end{corollary}

\begin{remark}\label{rem:X.1.5}
注意 \ref{thm:X.1.3} 的证明本质地用 \ref{prop:X.1.2}，从而用代数-形式几何的"第一比较定理". 相反，第 IX 讲的下降理论只通过 IX 3.4 出现，而在使 $f_*(\mathcal{O}_X) = \mathcal{O}_Y$ 的紧合态射 $f \colon X \to Y$ 的情形，\ref{cor:X.3.4} 的直接证明容易. 设 $Y'$ 平展于 $Y$ 上，假设 $X' = X \times_Y Y'$ 是两个非空开集的无交和，证明 $Y'$ 亦然. 确实有 $Y' = \mathrm{Spec}(\mathcal{A})$，故 $X' = \mathrm{Spec}(\mathcal{B})$，其中 $\mathcal{B} = \mathcal{A} \otimes_{\mathcal{O}_Y} \mathcal{O}_X$，而 $X'$ 的直和分解对应 $\mathcal{B}$ 分解成两个非零代数 $\mathcal{B}_1$ 和 $\mathcal{B}_2$ 的乘积. 由于 $f_*(\mathcal{O}_X) = \mathcal{O}_Y$，容易得出 $f_*(\mathcal{B}) = \mathcal{A}$，故 $\mathcal{A}$ 是两个代数 $f_*(\mathcal{B}_1)$ 和 $f_*(\mathcal{B}_2)$ 的和 (它们也非零，因为它们的单位截面非零)，证毕.
\end{remark}

\refstepcounter{theorem}\noindent\textbf{\thesection.\arabic{theorem}.}\quad 仍假设 $f$ 紧合且可分，但不再对 $f_*(\mathcal{O}_X)$ 作假设，它对应 $Y$ 的完全确定的平展覆盖 $Y'$，且在 $b$ 上方被 $a$ 的像 $b'$ 带点. 对典范态射 $X \to Y'$ 应用 \ref{cor:X.1.4}，并假设 $f$ 满射，正合序列 \ref{cor:X.1.4} 换成下列序列，它是代数拓扑中纤维空间的同伦正合序列的类似物：
\[
\pi_1(X_{\overline{y}}, \overline{a}) \longrightarrow \pi_1(X, a) \longrightarrow \pi_1(Y, b) \longrightarrow \pi_0(X_{\overline{y}}, \overline{a}) \longrightarrow \pi_0(X, a) \longrightarrow \pi_0(Y, b) \longrightarrow e.
\]
当然，在 \ref{cor:X.1.4} 中一般不能断言同态 $\pi_1(X_{\overline{y}}, \overline{a}) \to \pi_1(X, a)$ 单射；在代数拓扑中它的核是 $\pi_2(Y, b)$ 的像，在代数几何中也有必要引入各维数的同伦群，以及满足适当假设 (例如光滑态射) 的紧合态射的完整同伦正合序列. 目前除高同伦群的合理 (即使不是最终) 定义外，没有任何这方面的结果.

\begin{corollary}\label{cor:X.1.7}
设 $k$ 是代数闭域，$X$ 和 $Y$ 是 $k$ 上的两个连通预概形，假设 $X$ 在 $k$ 上紧合，$Y$ 局部诺特. 设 $a$ 是 $X$ 的几何点，$b$ 是 $Y$ 的在同一个 $k$ 的代数闭扩张 $K$ 中取值的几何点，考虑 $X \times_k Y$ 的几何点 $c = (a, b)$，以及由两个投影 $X \times_k Y \to X$ 和 $X \times_k Y \to Y$ 的基本群同态得到的同态 $\pi_1(X \times_k Y, c) \to \pi_1(X, a) \times \pi_1(Y, b)$. 前述同态是同构.
\end{corollary}

先假设 $K = k$. 令 $Z = X \times_k Y$，考虑投影 $f \colon Z \to Y$ 和 $Y$ 的几何点 $b$ 的局部化，对这种情况应用结果 1.4. 为此注意，换到 $X_{\mathrm{r\acute{e}d}}$ (不改变所考虑的基本群) 后可假设 $X$ 已既约故在 $k$ 上可分，故 $Z$ 在 $k$ 上可分，且显然有几何连通纤维 (因为 $X$ 连通). $Z$ 在 $b$ 处的几何纤维典范同构于 $X \otimes_k K = X$. 另一方面，由于复合态射 $X \to Z \to X$ 是恒等，得到 $\pi_1(X, a) \to \pi_1(Z, c)$ 单射，且 \ref{cor:X.1.4} 给出正合序列：
\[
e \longrightarrow \pi_1(X, a) \longrightarrow \pi_1(Z, c) \longrightarrow \pi_1(Y, b) \longrightarrow e.
\]
另一方面，有典范正合序列：
\[
e \longrightarrow \pi_1(X, a) \longrightarrow \pi_1(X, a) \times \pi_1(Y, b) \longrightarrow \pi_1(Y, b) \longrightarrow e,
\]
其中写的两个同态是典范单射和典范投影. 最后借助两端项上的恒等态射和中间项的典范同态 $\pi_1(Z, c) \to \pi_1(X, a) \times \pi_1(Y, b)$，有第一个正合序列到第二个的同态. 如此得到的图式的交换性平凡验证. 由于两端项上的同态是同构，中间项亦然，这在该情形证明 1.7.

当不再假设 $K = k$ 时，只得到同构
\[
\pi_1(Z, c) \longrightarrow \pi_1(X \otimes_k K, a) \times \pi_1(Y, b),
\]
而 \ref{cor:X.1.7} 等价于下列特殊情形：

\begin{corollary}\label{cor:X.1.8}
设 $X$ 是代数闭域 $k$ 上的紧合连通概形，$k'$ 是 $k$ 的代数闭扩张，$a'$ 是 $X \otimes_k k'$ 的几何点，$a$ 是它在 $X$ 中的像. 则典范同态 $\pi_1(X \otimes_k k', a') \to \pi_1(X, a)$ 是同构.
\end{corollary}

该同态满射的事实等价于：若 $X'$ 是 $X$ 的连通平展覆盖，则 $X' \otimes_k k'$ 也连通，它立即由 $k$ 代数闭的事实得出；它也是 IX 3.4 的特殊情形. 对 $X$ 的紧合性假设还没有用. 这样，说所考虑的同态单射也意味着：$X \otimes_k k'$ 的每个平展覆盖同构于 $X$ 的某个平展覆盖的逆像. 本质上是寻常的：可以找到 $k'$ 的有限型子 $k$-代数 $A$ 和 $X \otimes_k A$ 的平展覆盖，使它在 $X \otimes_k k'$ 上的逆像同构于给定覆盖. 设 $Y = \mathrm{Spec}(A)$，它是有限型的整 $k$-概形，故有 $k$ 上有理点. 对 $X \times Y$ 在 $k$ 上有理点 $(a, b)$ 处的基本群应用 \ref{cor:X.1.7}：得到 $X \times Y$ 的每个连通平展覆盖同构于覆盖 $X' \times Y'$ 的商，其中 $X'$，$Y'$ 分别是 $X$ 和 $Y$ 的群 $G$，$G'$ 的平展 Galois 覆盖，商按 $G \times G'$ 的子群 $H$. 这蕴含该 $X \times Y$ 覆盖在 $X \times Y'$ 上的逆像同构于形如 $X'_1 \times Y'$ 的覆盖，其中 $X'_1$ 是 $X$ 的平展覆盖. 故若 $L$ 是 $Y$ 的函数域，等于 $A$ 在 $k'$ 中的分式域，由给定 $X \times_k Y$ 覆盖诱导的 $X \otimes_k L$ 的平展覆盖，使存在 $L$ 的有限可分扩张 $L'$，使得所述覆盖在 $X \otimes_k L'$ 上的逆像同构于 $X'_1 \otimes_k L'$. 而 $k'$ 代数闭，可假设 $L$ 的扩张 $L'$ 包含在 $k'$ 中. 这证明给定的 $X \otimes_k k'$ 的平展覆盖同构于 $X'_1 \otimes_k k'$，证毕.

顺带指出的积 $X \times_k Y$ 的平展覆盖的显式形式立即蕴含下列结果：

\begin{corollary}\label{cor:X.1.9}
设 $k$ 是代数闭域，$X$ 和 $Y$ 是 $k$ 上两个局部诺特预概形，$Z = X \times_k Y$ 是它们的积，$Z'$ 是 $Z$ 的平展覆盖. 对每个在 $k$ 上有理的点 $y \in Y$，设 $i_y \colon \mathrm{Spec}(k) \to Y$ 是典范伴随的态射，$j_y = \mathrm{id}_X \times_k i_y$ 是对应的态射 $X \to Z$. 最后设 $X'_y$ 是 $Z'$ 经 $j_y$ 的逆像得到的 $X$ 的平展覆盖. 假设 $Y$ 连通，且 $X$ 或 $Y$ 在 $k$ 上紧合. 则 $X$ 的诸覆盖 $X'_y$ 全部同构.
\end{corollary}

形象地说，由连通预概形 $Y$ 参数化的 $X$ 的平展覆盖族是常的，若 $X$ 或参数预概形 $Y$ 在 $k$ 上紧合.

\begin{remark}\label{rem:X.1.10}
\cref{cor:X.1.7,cor:X.1.8,cor:X.1.9} 在正规代数概形的情形归功于 Lang-Serre \cite{exp10-2}. (他们的工作是这个讨论班发展的基本群理论的最初动机). 如这两位作者所指出的，放弃紧合性假设时这些结果变得不精确，至少在特征 $p > 0$ 时. 例如取 $X$ 为仿射直线 $X = \mathrm{Spec}(k[t])$，不难看到由方程
\[
x^p - x = st
\]
定义的、由仿射直线 $Y = \mathrm{Spec}(k[s])$ 参数化的 $X$ 的覆盖平展，且两两不同构. 这使 \ref{cor:X.1.9} 从而更不用说 \ref{cor:X.1.7} 失效，同样看到若 $s$ 看作 $k$ 的代数闭扩张 $K$ 中 $k$ 上的超越元，则找到 $X \otimes_k K$ 的平展覆盖 $X'$，它不来自 $X$ 的平展覆盖.
\end{remark}

\section{层存在定理的应用：紧合可分态射纤维基本群的半连续性定理}

\begin{theorem}\label{thm:X.2.1}
设 $Y$ 是剩余域 $k$ 的完备诺特局部环的谱，$X$ 是 $Y$-概形紧合，$X_0 = X \otimes_A k$，$a_0$ 是 $X_0$ 的几何点，$a$ 是 $X$ 的对应几何点. 则典范同态 $\pi_1(X_0, a_0) \to \pi_1(X, a)$ 是同构.
\end{theorem}

这只是在基本群语言中翻译 IX 1.10 中回顾的结果. 正是这里，代数-形式几何中的层存在定理本质地进入基本群理论.

现在引入剩余域 $k$ 的代数闭包 $\overline{k}$ 和几何纤维 $\overline{X}_0 = X_0 \otimes_k \overline{k}$. 有正合序列 (\cref{thm:IX.6.1})
\[
e \longrightarrow \pi_1(\overline{X}_0, \overline{a}) \longrightarrow \pi_1(X_0, a_0) \longrightarrow \pi_1(\overline{k}, \overline{k}) \longrightarrow e,
\]
另一方面有同构 \ref{thm:X.2.1} 和更初等的类似同构
\[
\pi_1(\overline{k}, \overline{k}) \longrightarrow \pi_1(Y, b),
\]
其中 $b$ 是 $a$ 在 $Y$ 中的像. 由此得到：

\begin{corollary}\label{cor:X.2.2}
用前面记号，假设 $\overline{X}_0 = X_0 \otimes_k \overline{k}$ 连通，设 $\overline{a}_0$ 是 $\overline{X}_0$ 的几何点，$a_0$ 是它在 $X$ 中的像，$b_0$ 是它在 $Y$ 中的像. 则下列典范同态序列
\[
e \longrightarrow \pi_1(\overline{X}_0, \overline{a}) \longrightarrow \pi_1(X, a_0) \longrightarrow \pi_1(Y, b_0) \longrightarrow e
\]
正合.
\end{corollary}

把这个序列与正合序列 \ref{cor:X.1.4} 比较，但注意 a) 对 $X \to Y$ 不必作平坦性或纤维可分性假设；b) 有重要补充：态射 $\pi_1(\overline{X}_0, \overline{a}_0) \to \pi_1(X, a_0)$ 单射.

这最后一个事实使我们能把 $X$ 在 $Y$ 上的其他几何纤维的基本群与 $\overline{X}_0$ 的比较. 设 $y_1$ 是 $Y$ 的任意点，$X_1$ 是它的纤维，$\overline{X}_1$ 是它的几何纤维 (相对 $k(y_1)$ 的某个代数闭扩张)，$a_1$ 是 $\overline{X}_1$ 的几何点，$a_1$ 是它在 $X$ 中的像，$b_1$ 是它在 $Y$ 中的像. 选择从 $a_1$ 到 $a_0$ 的"道路类"，从而有从 $b_1$ 到 $b_0$ 的道路类，从而有同态的交换图式：
\[
\begin{tikzcd}[row sep=large, column sep=large]
& \pi_1(\overline{X}_1, \overline{a}_1) \arrow[r] \arrow[d] & \pi_1(X, a_1) \arrow[r] \arrow[d] & \pi_1(Y, b_1) \arrow[r] & e \\
e \arrow[r] & \pi_1(\overline{X}_0, \overline{a}_0) \arrow[r] & \pi_1(X, a_0) \arrow[r] & \pi_1(Y, b_0) \arrow[r] & e,
\end{tikzcd}
\]
其中写的两个垂直箭头是同构. 由于第二行正合，找到典范同态，称之为基本群的特殊化同态\index[term]{jibenqunteshuhuatongtai@基本群特殊化同态} (只依赖所选的从 $a_1$ 到 $a_0$ 的道路类，故模 $\pi_1(X, a_0)$ 的内自同构确定)：
\[
\pi_1(\overline{X}_1, \overline{a}_1) \longrightarrow \pi_1(\overline{X}_0, \overline{a}_0).
\]
当上面第一行也正合时，立即得出特殊化同态满射. 因此考虑到 \ref{cor:X.1.4} 得到：

\begin{corollary}\label{cor:X.2.3}
在 \ref{thm:X.2.1} 的条件下，再假设态射 $f \colon X \to Y$ 可分 (\cref{def:X.1.1}) 且 $\overline{X}_0$ 连通 (故由 \ref{prop:X.1.2} 有 $f_*(\mathcal{O}_X) = \mathcal{O}_Y$). 则对 $X$ 在 $Y$ 上的每个几何纤维 $\overline{X}_1$ (配备几何点 $a_1$)，上面定义的特殊化同态是满同态.
\end{corollary}

这是基本群的半连续性结果，它在代数拓扑中似乎还没有类似物. 还可以在更一般的条件下陈述：

\begin{corollary}\label{cor:X.2.4}
设 $f \colon X \to Y$ 是纤维几何连通的紧合态射，$Y$ 局部诺特，$y_0$ 和 $y_1$ 是 $Y$ 的两个满足 $y_0 \in \overline{\{y_1\}}$ 的点，$\overline{X}_0$ 和 $\overline{X}_1$ 是对应 $k(y_0)$ 和 $k(y_1)$ 的给定代数闭扩张的几何纤维，$a_0$ resp. $a_1$ 是 $\overline{X}_0$ resp. $\overline{X}_1$ 的几何点. 则可以自然方式定义特殊化同态：
\[
\pi_1(\overline{X}_1, \overline{a}_1) \longrightarrow \pi_1(\overline{X}_0, \overline{a}_0),
\]
它在模内自同构意义下确定，且若 $f$ 是可分态射 (\cref{def:X.1.1}) 则是满同态.
\end{corollary}

事实上，首先由 \ref{cor:X.1.8} 得出 \ref{cor:X.2.4} 本质上与所选剩余域 $k(y_0)$ 和 $k(y_1)$ 的代数闭扩张无关. 这使我们能把 $Y$ 换成 $Y$ 上有位于 $y_0$ (resp. $y_1$) 上方的点 $y'_0$ (resp. $y'_1$) 的预概形 $Y'$. 取 $Y'$ 为 $Y$ 中 $y_0$ 的局部环的完备化的谱，应用 2.3.

\begin{remark}\label{rem:X.2.5}
\ref{cor:X.2.4} 关于特殊化同态满射的最终结论，更不用说作为其推论的结果 \ref{thm:X.1.3} 和 \ref{cor:X.1.4}，在不再假设 $f \colon X \to Y$ 是可分态射时变得不精确，即使在特征 0 的代数闭域上的射影概形情形. 我们后面将看到例子，既在 $f$ 平坦但 $f$ 有不可分纤维的情形 ($X$ 和 $Y$ 在 $k$ 上仍光滑)，也在 $f$ 的诸纤维确实可分但 $f$ 不平坦的情形 (例如 $f \colon X \to Y$ 是整正规概形的双有理态射)，参见 XI 3. 在这些例子中，一般几何纤维的基本群可能为零，而某个适当的特殊几何纤维的基本群非零. 另一方面，即使 $f \colon X \to Y$ 如 \ref{cor:X.2.4} 那样是紧合可分态射，特殊化态射也经常不是同构. 例如，容易给出 $\overline{X}_1$ 是非奇异椭圆曲线的例子 (故它的基本群交换，且对与特征互素的素数 $\ell$，它的 $\ell$-准素分量同构于 $\mathbb{Z}_\ell^2$，参见 XI)，而 $\overline{X}_0$ 由两条在两点相交的非奇异有理曲线组成，或由两条在一点相切的有理曲线组成，或最后由一条有尖点奇点的有理曲线组成. (关于退化椭圆曲线的完整分类，参见 Kodaira \cite{exp10-1} 和 Néron 的近期工作). 此时在这些情形中，$\overline{X}_0$ 的基本群分别是 $\widehat{\mathbb{Z}}$，$e$，$e$，故"严格小于"$\overline{X}_1$ 的. 我们后面将在 $f$ 是光滑态射时看到特殊化同态核的上界，它特别蕴含若 $k(y_0)$ 特征 0，则特殊化同态是同构. 但即使对光滑态射，若 $k(y_0)$ 的特征 $> 0$，特殊化同态也可能不是同构，例如 $X$ 是 $Y$ 上的阿贝尔概形 (若愿意，相对维数 1) 的情形，参见 XI 2. 基本群特殊化的满意理论必须考虑"真正的"基本群的"连续分量"，它对应无穷小群作为结构群的覆盖的分类；由此有理由期待：紧合光滑态射 $f \colon X \to Y$ 的几何纤维的"真正"基本群在 $X$ 上构成漂亮的局部系，它是 $X$ 上有限平坦群概形的投射极限\footnote{这个极具吸引力的猜想不幸被 M. Artin 的一个未发表例子否定，甚至在 $f$ 的诸纤维是给定亏格 $g \geqslant 2$ 的代数曲线时.}. 我们以后将回到这个观点，而当前的对象正相反，是尽量推进拓扑理论与基本群的概形理论共同的现象.
\end{remark}

现在设 $X_0$ 是代数闭域 $k$ 上亏格 $g$ 的紧合、光滑且连通的曲线. 若 $k$ 特征零，它的基本群可以用超越方式如下确定. 已知 $X_0$ 来自定义在素数域 $\mathbb{Q}$ 的有限超越次数的代数闭扩张上的曲线的基扩张，考虑到 \ref{cor:X.1.8}，可假设 $k$ 本身在 $\mathbb{Q}$ 上超越次数有限. 故可假设 $k$ 是复数域 $\mathbb{C}$ 的子域，再一次应用 \ref{cor:X.1.8} 可假设 $k = \mathbb{C}$. 此时不难验证 $X$ 的基本群同构于关联拓扑空间 $\tilde{X}$ (亏格 $g$ 的紧定向曲面) 的基本群关于有限指数子群定义的拓扑的完备化\footnote{这个推导在一个未整理的讲授中明确过.}. 另一方面，众所周知拓扑基本群由 $2g$ 个生成元 $s_i$，$t_i$ ($1 \leqslant i \leqslant g$) 生成，只服从一个关系：
\[
(s_1 t_1 s_1^{-1} t_1^{-1}) \cdots (s_g t_g s_g^{-1} t_g^{-1}) = 1.
\]
故 $X$ 的基本群有 $2g$ 个拓扑生成元 $s_i$，$t_i$ ($1 \leqslant i \leqslant g$)，只服从前面关系. 若 $k$ 的特征为 $p > 0$，用 $A$ 记由 $k$ 构造的 Witt 向量环，用 $K$ 记它的分式域的代数闭扩张. 在 III 7.4 中已见存在 $Y = \mathrm{Spec}(A)$ 上的概形 $X$，在 $Y$ 上紧合且光滑，约化为 $X_0$. 对它应用 \ref{cor:X.2.3}，找到满态射
\[
\pi_1(X_1) \longrightarrow \pi_1(X_0),
\]
其中 $X_1 = X \otimes_A K$. 立即\footnote{参见 EGA IV 12.2.} $X_1$ 在 $K$ 上光滑、连通 (\cref{prop:X.1.2})、维数 1，且它的亏格等于 $g$ (由 Euler-Poincaré 特征的不变性，参见 EGA III 7). 由于 $K$ 特征 0，可以应用前面结果. 于是用超越方式证明：

\begin{theorem}\label{thm:X.2.6}
设 $X_0$ 是代数闭域 $k$ 上的光滑紧合连通代数曲线，$g$ 是它的亏格. 则 $\pi_1(X_0)$ 有 $2g$ 个拓扑生成元的生成系，只服从上面写的关系. 当 $k$ 的特征为 0 时，$\pi_1(X_0)$ 甚至是前面生成元和关系的自由 Galois 型群.
\end{theorem}

\begin{remark}\label{rem:X.2.7}
据编写者所知，目前没有用纯代数方式证明前面结果的方法 (除亏格 0，1 外). 首先，看不出如何在 $\pi_1(X_0)$ 中区分 $2g$ 个元素，然后才能期望它们构成拓扑生成元系. 在这方面，去掉 $n$ 个点的有理直线及其中在这 $n$ 点温和分歧的覆盖研究的情形更友好，因为考虑在这 $n$ 点的分歧群提供待研究基本群的 $n$ 个元素，确实证明它们拓扑生成该基本群，如我们后面将看到的\footnote{参见第十二讲. 而且这些元素只在模共轭意义下真正确定，需要在它们的类中明智地同时选择.}. 但即使在这个特别具体的情形，似乎也不存在纯代数证明. 这样的证明显然极其有趣. 关于曲线基本群唯一能用纯代数方式证明的事实 (除下面 \ref{cor:X.2.12} 的弱有限性定理，Lang-Serre \cite{exp10-2} 用代数方式证明)，似乎是通过 Jacobi 簇确定化为阿贝尔群的基本群 (在 IX 5.8 末行指出).
\end{remark}

\refstepcounter{theorem}\noindent\textbf{\thesection.\arabic{theorem}.}\quad \ref{thm:X.2.6} 的最后一个断言在特征 $p > 0$ 时不再成立，如在椭圆曲线情形已经看到. 如已指出的，我们不知道 $X_0$ 的基本群是否拓扑有限表现；这似乎完全不可能.

\begin{theorem}\label{thm:X.2.9}
设 $k$ 是代数闭域，$X$ 是 $k$ 上紧合连通概形. 则 $X$ 的基本群拓扑有限生成.
\end{theorem}

对 $n = \dim X$ 作递归，$n \leqslant 0$ 时断言平凡. 假设 $n > 0$ 且定理对维数 $n' < n$ 已证. 由 Chow 引理 (EGA II 5.6.2) 存在 $k$ 上射影概形 $X'$ 和满射态射 $X' \to X$. 显然可假设 $X'$ 既约，再换到正规化后正规. 借助下降理论，只需证明 $X'$ 的诸连通分量的基本群拓扑有限生成 (\cref{cor:IX.5.2}). 这归结到 $X$ 射影且正规的情形. 若 $n = 1$，只需应用 2.6. 若 $n \geqslant 2$，考虑射影浸入 $X \to \mathbf{P}^r_k$，和满足 $Y \neq X$ (i.e. $H \not\supset X$) 的超平面截面 $Y = X \cdot H$ (配备诱导既约结构). 则 $\dim Y < n$，考虑到递归假设，只需证明 $\pi_1(Y) \to \pi_1(X)$ 满射. 更一般地：

\begin{lemma}\label{lem:X.2.10}
设 $X$ 是代数闭域 $k$ 上的紧合预概形，$g \colon X \to \mathbf{P}^r_k$ 是态射. 假设 $X$ 不可约且正规，$\dim g(X) \geqslant 2$. 设 $H$ 是 $\mathbf{P}^r_k$ 的超平面，$Y = X \times_{\mathbf{P}^r_k} H$. 则 $Y$ 连通，且同态 $\pi_1(Y) \to \pi_1(X)$ 满射.
\end{lemma}

这些断言确实由下列断言得出：

\begin{corollary}\label{cor:X.2.11}
在先前条件下，设 $X'$ 是 $X$ 的连通平展覆盖，$Y' = X' \times_X Y = X' \times_{\mathbf{P}^r_k} H$ 是 $Y$ 上诱导的覆盖. 则 $Y'$ 连通.
\end{corollary}

由于 $X$ 正规，$X'$ 正规，故连通时 $X'$ 不可约；且它在 $\mathbf{P}^r_k$ 中的像维数 $\geqslant 2$. 一个归功于 Zariski 的熟知引理 (称为"Bertini 定理") 蕴含：若 $H_1$ 是 $\mathbf{P}^r_k$ 中定义在 $k$ 的扩张 $K$ 上的一般超平面，则 $X' \times_{\mathbf{P}^r} H_1$ 在 $K$ 上万有不可约，故万有连通. Zariski 连通定理 (EGA III 4) 蕴含对每个超平面 $H$ (定义在 $k$ 的任意扩张上) $X' \times_{\mathbf{P}^r} H$ 几何连通. 这完成 \ref{cor:X.2.11} 的证明，从而 \ref{thm:X.2.9} 的证明.

\begin{corollary}\label{cor:X.2.12}
(Lang-Serre) 在 \ref{thm:X.2.9} 的条件下，对每个有限群 $G$，$X$ 的群 $G$ 的主覆盖的同构类的集有限.
\end{corollary}

\begin{remark}\label{rem:X.2.13}
在 \ref{lem:X.2.10} 的条件下，当 $\dim g(X) \geqslant 3$ (至少在 $g$ 是浸入且 $X$ 正则时) 我们将证明更精确的结果，在代数几何中以"Lefschetz 定理"之名著称：$\pi_1(Y) \to \pi_1(X)$ 是同构.\footnote{参见本讨论班的续篇 SGA 2 (1962)，定理 X 3.10.} 在经典情形对同调群和高同伦群有类似断言，它们迟早应被纳入抽象代数几何. 即使对 Hodge 上同调 $H^p(X, \Omega^q)$，问题似乎也还未被研究；对后者，Lefschetz 定理在特征 $p > 0$ 时原样成立也不太可能.

\begin{remark}[2003 年 (MR) 增补]
设 $R$ 是剩余域 $k$ 代数闭且特征 $p > 0$、分式域 $K$ 的完备离散赋值环，设 $Y$ 是 $R$ 上亏格 $g$ 的紧合光滑连通曲线. 有满的特殊化态射 $sp \colon \pi_1(Y_K) \to \pi_1(Y_k)$. 已注意到若 $K$ 特征 0，则只要 $g \geqslant 1$，$sp$ 就不是同构. 假设 $K$ 特征 $p$，故 $R$ 同构于形式幂级数环 $k[\![T]\!]$. 在等特征 $p > 0$ 时，基本群不由亏格 $g$ 确定，如在椭圆曲线 (可以是普通或超奇异的) 上已经看到. 引用 A. Tamagawa 的近期结果 (尚未发表). 若 $G$ 是 pro-有限群，记 $G^{\mathrm{r\acute{e}s}}$ 为 $G$ 的 pro-有限商群，即 $G$ 的有限可解拓扑商的投射极限.

\begin{theorem}[A. Tamagawa]
假设 $g \geqslant 2$，特殊纤维 $Y_k$ 在有限域上可定义，且由 $sp$ 取商得到的态射 $sp^{\mathrm{r\acute{e}s}} \colon \pi_1(Y_K)^{\mathrm{r\acute{e}s}} \to \pi_1(Y_k)^{\mathrm{r\acute{e}s}}$ 双射. 则曲线 $Y$ 在 $R$ 上为常.
\end{theorem}

注意 $\overline{K}/K$ 的 伽罗瓦群可解. 前述断言还可以翻译成：假设一般纤维 $Y_K$ 的每个有限平展覆盖 $X_K \to Y_K$ (可解 伽罗瓦群的 Galois 覆盖) 都有潜在好约化 (即在 $K$ 的某个有限扩张中用 $R$ 在其中的正规化替换 $R$ 后，延拓成 $Y$ 的有限平展覆盖). 则 $Y$ 在 $R$ 上为常.
\end{remark}
\end{remark}

\section{纯度定理的应用：紧合光滑态射纤维基本群的连续性定理}

不证地回忆\footnote{证明参见 SGA 2 X 3.4.}：

\begin{theorem}\label{thm:X.3.1}
[\emph{纯度定理} (Zariski-Nagata)]\index[term]{chundudingli@纯度定理} 设 $f \colon X \to Y$ 是整预概形的拟有限支配态射，$X$ 正规，$Y$ 正则局部诺特，设 $Z$ 是 $f$ 不平展 (i.e. $f$ 分歧，两者相同，I 9.5 (ii)) 的 $X$ 的点集. 若 $Z \neq X$，则 $Z$ 在它的每个点处余维数 1，i.e. 对 $Z$ 的每个不可约分量 $Z'$ (生成点 $z$)，$\mathcal{O}_{X,z}$ 的 Krull 维数等于 1.
\end{theorem}

回忆预概形称为正规 resp. 正则，若它的局部环正规 resp. 正则，且关系 $Z \neq X$ 也意味着有限扩张 $R(Z)/R(X)$ (其中 $R$ 表示有理函数域) 可分. 在 $Z$ 的分量 $Z'$ 的生成点 $z$ 处，并在 $z$ 下方的 $Y$ 的点 $y$ 处局部化，得到等价断言：

\begin{corollary}\label{cor:X.3.2}
设 $A$ 是诺特正则局部环，$A \to B$ 是单射局部同态，使 $B$ 正规，是 $A$ 上有限型代数的局部化，且在 $A$ 上拟有限；再假设 $\dim A (= \dim B) \geqslant 2$，且对 $B$ 的每个异于极大理想的素理想 $\mathfrak{p}$，$B$ 在 $\mathfrak{p}$ 处平展于 $A$ 上，i.e. $B_{\mathfrak{p}}$ 在 $A_{\mathfrak{q}}$ (其中 $\mathfrak{q} = A \cap \mathfrak{p}$) 上平展. 则 $B$ 在 $A$ 上平展.
\end{corollary}

而且不难把最后一个断言归结到 $A$ 是完备局部环、故 $B$ 在 $A$ 上有限的情形. Zariski \cite{exp10-5} 给出这个结果的简单证明，在等特征情形有效；一般情形归功于 Nagata \cite{exp10-3}，他依靠 Chow 的精细结果；后者没有被讨论班的任何参加者验证过，应成为以后讲授的对象. 这里只指出在 $\dim A = 2$ 的特殊情形下非常简单的证明，它足以用于我们本节的这个最重要应用. 由于 $B$ 正规，它是深度 (旧术语：上同调余维数) $\geqslant 2$ 的 $B$-模，故也是深度 $\geqslant 2$ 的 $A$-模，且 $A$ 维数 2 正则，由此得出 $B$ 是 $A$ 上的自由模\footnote{参见 EGA IV$_0$ 17.3.4.}. 于是由 I 4.10 得出 $B$ 在 $A$ 上分歧的素理想 $\mathfrak{q}$ 的集是 $\mathrm{Spec}(A)$ 的由主理想 (由 $B$ 在 $A$ 上的基的判别式生成) 定义的部分，故它包含在 $\mathrm{Spec}(A)$ 的闭点中时为空，这证明 $\dim A = 2$ 时的 3.2.

主要用 \ref{thm:X.3.1} 的等价形式：

\begin{corollary}\label{cor:X.3.3}
设 $X$ 是局部诺特正则预概形，$U$ 是 $X$ 的余维数 $\geqslant 2$ 的闭部分 $Z$ 的余集的开部分. 则函子 $X' \mapsto X' \times_X U$ (从 $X$ 的平展覆盖范畴到 $U$ 的平展覆盖范畴) 是范畴等价；特别，若 $a$ 是 $U$ 的几何点，典范同态 $\pi_1(U, a) \to \pi_1(X, a)$ 是同构.
\end{corollary}

最后一个断言显然是第一个的推论，对第一个显然可假设 $X$ 连通故不可约. 由 $X$ 的正规性已得函子 $X' \mapsto X' \times_X U$ (从 $X$ 的局部自由 (未必平展) 覆盖范畴到 $U$ 的覆盖范畴) 忠实全，因为函子 $E \mapsto E|_U$ (从 $X$ 上局部自由模范畴到 $U$ 上局部自由模范畴) 忠实全. 剩下证明对 $U$ 的每个平展覆盖 $U'$，存在 $X$ 的平展覆盖 $X'$ (由前面所述必然唯一)，使 $U'$ 同构于 $X' \times_X U$. 显然可假设 $U'$ 连通，故不可约，因为 ($U$ 正规时) $U'$ 正规. 设 $K$ 是 $X$ (或 $U$，相同) 上的有理函数域，$K'$ 是 $U'$ 的，则 $U'$ 等同于 $U$ 在 $K'$ 中的正规化 (\cref{cor:I.10.3}). 设 $X'$ 是 $X$ 在 $K'$ 中的正规化 (EGA II 6.3)，则 $X' \times_X U \xrightarrow{\sim} U'$，另一方面 $X'$ 正规、整，结构态射 $f \colon X' \to X$ 有限且支配 (因为 $X$ 正规且 $K'/K$ 是有限可分扩张). 它在 $f^{-1}(U) = X' - f^{-1}(Z)$ 上平展，且由于 $Z$ 在 $X$ 中余维数 $\geqslant 2$，$f^{-1}(Z)$ 在 $X'$ 中余维数 $\geqslant 2$. 于是由 \ref{thm:X.3.1} 得出 $X'$ 在 $X$ 上平展，这完成证明.

现在设 $f \colon X \to Y$ 是从局部诺特正则预概形 $X$ 到预概形 $Y$ 的有理映射，假设 $f$ 在余维数 $\geqslant 2$ 的闭部分的余集的开集 $U$ 上定义. 则由 \ref{cor:X.3.3} 推出函子 (在同构意义下确定) 从 $Y$ 的平展覆盖范畴到 $X$ 的平展覆盖范畴，从而对 $U$ 的每个像为 $b$ 的几何点 $a$，有典范同态
\[
\pi_1(X, a) \longrightarrow \pi_1(Y, b)
\]
(由典范同态 $\pi_1(U, a) \to \pi_1(Y, b)$ 借助同构 $\pi_1(U, a) \xrightarrow{\sim} \pi_1(X, a)$ 得到). 当 $f$ 是支配态射，$X$ 和 $Y$ 整、函数域为 $K$ 和 $L$ (故 $K$ 是 $L$ 的扩张)，$Y$ 正规时，这些对应借助域扩张明确：注意对 $L$ 的每个在 $Y$ 上非分歧的有限扩张 $L'$，代数 $K' = L' \otimes_L K$ 在 $X$ 上非分歧.

特别，这些思考表明：局部诺特连通正则预概形 (由余维数 $\leqslant 1$ 处局部化的几何点带点) 的基本群，在取该范畴中的态射为在余维数 $\geqslant 2$ 闭部分的余集上定义的支配有理映射时，是函子. 回忆例如域 $k$ 上正规概形到 $k$ 上紧合概形的有理映射在余维数 $\geqslant 2$ 的集的余集上定义，得到：

\begin{corollary}\label{cor:X.3.4}
(基本群的双有理不变性)\index[term]{jibenqunshuangyoulibubianxing@基本群双有理不变性} 设 $k$ 是域，$X$ 和 $Y$ 是 $k$ 上紧合的正则概形，$f \colon X \to Y$ 是从 $X$ 到 $Y$ 的双有理映射，$\Omega$ 是 $X$ 的函数域 $K$ 的代数闭扩张，用来定义 $X$ 的基本群和 $Y$ 的基本群. 则它们典范同构.
\end{corollary}

这也意味着：对 $K$ 的有限扩张 $K'$，若它在 $K$ 的某个紧合非奇异"模型" $X$ 上非分歧，则它在任何其他紧合非奇异模型上都非分歧.

\begin{remark}\label{rem:X.3.5}
关于纯度定理的其他应用，参见 Abhyankar 的工作，在 \cite{exp10-4} 中阐述，受 Zariski \cite[第八章]{exp10-6} 的结果启发，后者用拓扑方式证明. 它们还远未被"抽象"代数几何同化，值得新的努力.
\end{remark}

将需要分歧理论的一些初等事实. 设 $V$ 是分式域 $K$、剩余域 $k$ 的离散赋值环，$L$ 是 $K$ 的群 $G$ 的 伽罗瓦扩张，$V'$ 是 $V$ 在 $L$ 中的正规化 (它是 $V$ 上秩 $n = [L : K]$ 的自由模)，$\mathfrak{m}'$ 是 $V'$ 的极大理想，$G_d$ 是 $G$ 的保持 $\mathfrak{m}'$ 不变的元素的子群，故 $G_d$ 作用在 $k$ 的剩余扩张 $k' = V'/\mathfrak{m}'$ 上，$G_i$ 是 $G_d$ 的平凡作用的元素的子群 (回忆 $G_d$ 和 $G_i$ 分别称为 $G$ 的分解子群和惯性子群). 称 $L$ 在 $V$ 上温和分歧，若 $n_i = [G_i : e]$ 的阶与 $k$ 的特征互素 (若 $k$ 特征 0 该条件总满足). 众所周知此时 $G_i$ 典范嵌入 $k'^{*}$，故同构于 $k'$ 中 $n_i$ 次单位根群，这特别蕴含 $G_i$ 循环. 这种情况的典型情形是令 $L = K[t]/(t^n - u)$，其中 $u$ 是 $V$ 的幺模化元，$n$ 是与 $p$ 互素的整数：若 $K$ 含 $n$ 次单位根，$L$ 是 $K$ 的完全分歧 伽罗瓦扩张，伽罗瓦群 $G = G_i$ 同构于 $\mathbb{Z}/n\mathbb{Z}$.

\begin{lemma}\label{lem:X.3.6}
(Abhyankar 引理) 设 $V$ 是分式域 $K$ 的离散赋值环，$L$ 和 $K'$ 是 $K$ 的在 $V$ 上温和分歧的两个 伽罗瓦扩张，$n$ 和 $m$ 是对应惯性群的阶，$L'$ 是 $L$ 和 $K'$ 在 $K$ 上的复合扩张. 若 $m$ 是 $n$ 的倍数，则 $L'$ 在 $K'$ 的正规化 $V'$ 的诸局部化上非分歧.
\end{lemma}

设 $W'$ 是 $V'$ 在 $L'$ 中的正规化，$\mathfrak{m}'$ 是 $V'$ 的极大理想，$\mathfrak{n}'$ 是 $W'$ 的位于 $\mathfrak{m}'$ 上方的极大理想，$\mathfrak{n}$ 是它在 $V$ 在 $L$ 中的正规化 $W$ 上诱导的极大理想，$G$，$H$，$M$ 是 $L$，$K'$，$L'$ 在 $K$ 上的 伽罗瓦群，$G_i$，$H_i$，$L'_i$ 是对应所选极大理想的惯性群. 则 $M$ 嵌入乘积 $G \times H$，$M_i$ 嵌入乘积 $G_i \times H_i$，且投影 $M \to G$ 和 $M \to H$，$M_i \to G_i$ 和 $M_i \to H_i$ 满射 (中间域推论). 由此已经得出，由于由假设 $G_i$ 和 $H_i$ 是阶 $m$ 和 $n$ (与 $p$ 互素) 的循环群，$M_i$ 的阶与 $p$ 互素，故循环，且由于 $m$ 是 $n$ 的倍数故 $G_i \times H_i$ 的元素 $m$ 次幂为零，$M_i$ 的阶整除 $m$，故等于 $m$，因为 $M_i \to H_i$ 满射. 故这最后一个同态也单射. 而它的核是 $\mathfrak{n}'$ 在 $\mathfrak{m}'$ 上的惯性群，这证明 $L'$ 在 $\mathfrak{n}'$ 处非分歧于 $K'$ 上. 由此得引理.

现在置于 \ref{cor:X.2.4} 的条件下，有特殊化同态
\[
\pi_1(\overline{X}_1, \overline{a}_1) \longrightarrow \pi_1(\overline{X}_0, \overline{a}_0)
\]
它相对紧合可分态射 $f \colon X \to Y$ 满射. 想明确这个同态的核. 如 \ref{cor:X.2.4} 的证明那样进行，看到在这个问题中总可以假设 $Y$ 是离散赋值环 $V$ 的谱，完备且剩余域代数闭 (因为总可以找到这样的环和它的谱到 $Y$ 的、像为 $\{y_0, y_1\}$ 的态射). 则 $X_0 = \overline{X}_0$，$k(y_0) = k = V$ 的剩余域，$k(y_1) = K = V$ 的分式域. 设 $K_s$ 是 $K$ 的可分闭包，$\overline{K}$ 是它的代数闭包，对 $K$ 的每个含 $V$ 的子环 $W$，令 $X_W = X \otimes_V W$. 特别有
\[
X_V = X, \quad X_K = X_1, \quad X_{\overline{K}} = \overline{X}_1.
\]
而且典范态射 $\overline{X}_1 = X_{\overline{K}} \to X_{K_s}$ 诱导基本群同构 (\cref{cor:IX.4.11})，故考虑到同构 \ref{thm:X.2.1} $\pi_1(X_0) \to \pi_1(X)$，归结到研究与典范态射 $X_{K_s} \to X$ 关联的满同态
\[
\pi_1(X_{K_s}) \longrightarrow \pi_1(X).
\]
确定后者的核归结到下列问题的解：有 $X_{K_s}$ 的群 $G$ 的连通主覆盖 $Z_{K_s}$ (故对应 $\pi_1(X_{K_s})$ 到 $G$ 的同态)，确定在什么条件下它同构于 $X$ 的群 $G$ 的主覆盖 $Z$ 的原像.

先注意 $K_s$ 是它的有限子扩张 $K'$ 在 $K$ 上的递增滤过并，从而 $Z_{K_s}$ 同构于 $X_{K'}$ 的某个主覆盖 $Z_{K'}$ (对适当的 $K'$) 的逆像 (注意对固定的 $K'$，$Z_{K'}$ 不唯一确定). 说 $Z_{K_s}$ 同构于 $X$ 的主覆盖 $Z$ 的逆像，意味着存在 $K_s$ 的有限子扩张 $K'' \supset K'$，使 $Z_{K''} = Z_{K'} \otimes_{K'} K''$ 同构于 $Z \otimes_V K''$.

现在对 $K_s$ 的有限子扩张 $K'$，用 $V'$ 记 $V$ 在 $K'$ 中的正规化，它是离散赋值环，完备，剩余域 $k$. 故典范态射 $X_{V'} \to X_V$ 对 $Y = \mathrm{Spec}(V)$ 和 $Y' = \mathrm{Spec}(V')$ 的闭点上方的纤维诱导同构，于是由应用于 $X_V$ 和 $X_{V'}$ 的 \ref{thm:X.2.1} 得出基本群诱导同态 $\pi_1(X_{V'}) \to \pi_1(X_V)$ 是同构，或也即：$X_{V'}$ 的每个主覆盖是同构意义下确定的 $X_V$ 的主覆盖的逆像. 这蕴含

\begin{lemma}\label{lem:X.3.7}
设 $Z_{K'}$ 是 $X_{K'}$ 的群 $G$ 的连通主覆盖，$Z_{K_s}$ 是它在 $X_{K_s}$ 上的逆像. 为使后者同构于 $X$ 的主覆盖 $Z$ 的逆像，必要且充分条件是存在 $K$ 在 $K_s$ 中的有限扩张 $K'' \supset K'$，使 $X_{K''}$ 的主覆盖 $Z_{K''}$ 由 $X_{V''}$ 的主覆盖诱导.
\end{lemma}

特别假设诸 $X_{V''}$ 正规 (例如只需 $X_0$ 正规，更不用说 $X_0$ 简单，参见 I 9.1). 由于它们连通，故不可约. 设 $L$ 是 $X$ 和 $X_K$ 的有理函数域，$L'$ 是 $X_{V'}$ 和 $X_{K'}$ 的，$L''$ 是 $X_{V''}$ 和 $X_{K''}$ 的. 则在 \ref{lem:X.3.7} 的条件下，$Z_{K'}$ 定义 $L'$ 的有限可分扩张 $R'$，$Z_{K''}$ 定义扩张 $R'' = R' \otimes_{L'} L'' = R' \otimes_{K'} K''$. \ref{lem:X.3.7} 中考虑的条件于是也意味着存在 $K'$ 的有限可分扩张 $K''$，使 $R'' = R' \otimes_{K'} K''$ 在函数域 $L'' = L' \otimes_{K'} K''$ 的正规概形 $X_{V''}$ 上非分歧，而不只是在 $X_{V''}$ 的开部分 $X_{K''}$ 上非分歧.

此后假设 $f \colon X \to Y$ 是光滑态射，故诸态射 $X_{V'} \to \mathrm{Spec}(V')$ 光滑，故诸概形 $X_{V'}$ 正则. 注意 $\mathrm{Spec}(V')$ 的闭点在 $X_{V'}$ 中的纤维不可约且余维数 1. 设 $o'$ 是它的局部环，它是函数域 $L'$ 的离散赋值环，剩余域同构于 $X_0$ 的有理函数域，故与 $k$ 有同样特征. 类似地在 $L''$ 中定义 $o''$，它显然是 $o'$ 在 $L''$ 中的正规化. 于是由纯度\cref{thm:X.3.1} 或 \ref{cor:X.3.3} 得出：为使 $R''$ 在 $X_{V''}$ 上非分歧，必要且充分条件是 $R''$ 在 $o''$ ($o'$ 在 $L''$ 中的正规化) 上非分歧.

现在注意，若 $u'$ 是 $V'$ 的幺模化元，它也是 $o'$ 的幺模化元. 若 $n$ 是与 $k$ 的特征 $p$ 互素的整数，取 $K'' = K'[t]/(t^n - u)$，则 $K''$ 是 $K'$ 的有限 伽罗瓦扩张，且 $L''$ 同构于 $L'[t]/(t^n - u')$，故在 $o'$ 上温和分歧，惯性群阶 $n$. 假设 $G$ 的阶与 $p$ 互素，这蕴含 $R'$ 在 $o'$ 上温和分歧，取 $n$ 为 $R'$ 在 $o'$ 上的惯性群阶的与 $p$ 互素的倍数 (例如 $n = [G : e]$). 应用 Abhyankar \cref{lem:X.3.6}，看到 \ref{lem:X.3.7} 中考虑的条件满足.

这证明下列定理：

\begin{theorem}\label{thm:X.3.8}
设 $f \colon X \to Y$ 是纤维几何连通的紧合光滑态射，$Y$ 局部诺特，$y_0$ 和 $y_1$ 是满足 $y_0 \in \overline{\{y_1\}}$ 的 $Y$ 的两个点，$\overline{X}_0$ 和 $\overline{X}_1$ 是对应的几何纤维，考虑特殊化同态 (\cref{cor:X.2.4}) $\pi_1(\overline{X}_1) \to \pi_1(\overline{X}_0)$. 该同态满射，且 $\pi_1(\overline{X}_1)$ 到阶与 $k(y_0)$ 的特征 $p$ 互素的有限群 $G$ 的每个连续同态都来自 $\pi_1(\overline{X}_0)$ 到 $G$ 的同态.
\end{theorem}

换言之：

\begin{corollary}\label{cor:X.3.9}
若 $k(y_0)$ 特征零，则特殊化同态是同构. 若 $k(y_0)$ 特征 $p > 0$，则特殊化同态的核包含在 $\pi_1(\overline{X}_1)$ 到阶与 $p$ 互素的有限群的连续同态的核的交中 (或即：Galois 型群 $\pi_1(\overline{X}_1)$ 的 Sylow $p$-子群生成的闭不变子群)；故若 $\pi_1(\overline{X}_1)^{(p)}$ 表示 $\pi_1(\overline{X}_1)$ 除以前述闭子群的商群，且类似地定义 $\pi_1(\overline{X}_0)^{(p)}$，则特殊化同态诱导同构
\[
\pi_1(\overline{X}_1)^{(p)} \xrightarrow{\sim} \pi_1(\overline{X}_0)^{(p)}.
\]
\end{corollary}

注意 \ref{thm:X.3.8} 的证明纯代数. 如 \ref{thm:X.2.6} 那样进行，用超越方式由此得出：

\begin{corollary}\label{cor:X.3.10}
设 $X_0$ 是特征 $p$ 的代数闭域上亏格 $g$ 的紧合光滑连通曲线. 用 \ref{cor:X.3.9} 引入的记号，群 $\pi_1(X_0)^{(p)}$ 同构于 $\Gamma^{(p)}$，其中 $\Gamma$ 是由生成元 $s_i$，$t_i$ ($1 \leqslant i \leqslant g$) 服从关系
\[
(s_1 t_1 s_1^{-1} t_1^{-1}) \cdots (s_g t_g s_g^{-1} t_g^{-1}) = 1
\]
生成的 Galois 型群.
\end{corollary}

\begin{remark}\label{rem:X.3.11}
在 $k(y_0)$ 特征零的情形，结果 \ref{cor:X.3.9} 用超越方式众所周知. 注意 \ref{cor:X.3.10} 的证明在不等特征情形调用纯度定理，但只在维数 2 的环的情形，其中该定理的证明容易且在正文中已回忆.
\end{remark}

\section{文献}\printbibliography[heading=none]


